\section{\texorpdfstring{Structure of higher-genus GW theory of local \(\P^1\)}{Structure of higher-genus GW theory of local P1} (Shuai Guo)}%
\label{sec:Structure of higher-genus GW theory of local P1 (Shuai Guo)}

This lecture is based on joint work with Jingyi Xu and Qingsheng Zhang.

\subsection{The geometry of \(X_p\)}%
\label{sub:The geometry of Xp}

We will consider the total space
\[ X_p \coloneqq \on{Tot}\ab(\mc{O}_{\P^1}(p-1) \oplus \mc{O}_{\P^1}(-p-1)), \]
which is a non-compact toric Calabi-Yau threefold. For \(p = 0\) this is the resolved conifold.

Let \(N \simeq \Z^3\) be a lattice of rank \(3\) and let \(M \coloneqq \Hom(N,\Z)\) be the dual lattice, which we identify with the characters of \(\T \simeq (\C^{\times})^3\). The one-dimensional cones of the fan of \(X_p\) are generated by
\[ b_1 = (1-p,-1,1), \qquad b_2 = (0,1,1), \qquad b_3 = (1,0,1), \qquad b_4 = (0,0,1), \]
and \(X_p\) contains \(\T\) as an open dense subset. Using the dual vector \(e_3^{\vee} \in M\) we obtain the \textit{Calabi-Yau subtorus}
\[ \T' \coloneqq \ker\ab(\chi^{e_3^{\vee}} \colon \T \to \C^{\times}) \simeq (\C^{\times})^2, \]
whose action preserves the Calabi-Yau form. The toric graph consists of two triangles, which correspond to the two fixed points of the torus action. The weights of the \(\T'\)-action at the first fixed point are
\[ (0,1), \qquad (1,0), \qquad (-1,-1), \]
and at the second fixed point they are
\[ (0,-1), \qquad (1,1-p), \qquad (-1,p). \]
We can therefore express the invariants in terms of equivariant cohomology by localizing at these two fixed points. Here, \(H^*_{\T'}(\mr{pt}) = \C[\lambda_1,\lambda_2]\).

Unfortunately, for \(p \geq 2\) the toric graph of \(X_p\) is not convex. Then \(X_p\) is not semiprojective. This is the reason the case at hand is not covered by the existing literature, as we will explain in~\Cref{sub:Topological recursion and the remodeling conjecture}.

\subsection{The conjecture of CGMPS}%
\label{sub:The conjecture of CGMPS}

In 2007, Caporaso, Griguolo, Mari\~{n}o, Pasquetti, and Seminara studied the Gromov-Witten theory of \(X_p\). Writing
\[ \check{F}^{X_p}(w,\hbar) \coloneqq \sum_{g \geq 0} \check{F}_g^{X_p}(w) \hbar^{2g-2} \]
for the total free energy, they computed the genus-one part exactly and proposed the ansatz
\begin{equation}\label{eqn:cgmps ansatz}
    \check{F}_g^{X_p} = \frac{\wt{\mc{P}}_g(w,p)}{(w-w_c)^{5g-5}}, \qquad \wt{\mc{P}}_g(w,p) = \sum_{i=1}^{5g-5} a_{g,i}(p)(w-1)^i
\end{equation}
for \(g \geq 2\). Here, the variable \(w\) is related to the Kähler parameter \(q\) by \(w = 1-q\) and \(w_c = 1 - \frac{1}{p^2}\). Our convention will differ by \(F_g^{X_p} = (-1)^{g-1}\check{F}_g^{X_p}\) for \(g \geq 2\).

Analyzing the asymptotics of the genus-zero invariants, CGMPS defined the scaled string coupling by
\[ z^{5/2} = \hbar^{-2} \frac{p^8}{4w_c^5}(w-w_c)^5 \]
and introduced the \textit{double-scaling limit} \(F_{\mr{ds}}(z)\), which is obtained from \(\check{F}^{X_p}(w,\hbar)\) by letting \(w \to w_c\) and \(\hbar \to 0\) with \(z\) fixed. Combining their genus-zero and genus-one results with the ansatz gives
\[ F_{\mr{ds}}(z) = -\frac{4}{15}z^{5/2} - \frac{1}{48}\log z + \sum_{g \geq 2} b_g z^{-5(g-1)/2}, \]
where \(b_g = \ab(\frac{p^8}{4w_c^5})^{g-1}\wt{\mc{P}}_g(w_c,p)\).

\begin{conj}[CGMPS {\cite{CGMPS2007}}]%
\label{conj:cgmps}
    The double-scaling limit \(F_{\mr{ds}}(z)\) coincides with the free energy \(F_{(3,2)}(z)\) of the \((3,2)\) minimal model.
\end{conj}

The free energy of the \((3,2)\) minimal model has the form
\[ F_{(3,2)}(z) = -\frac{4}{15}z^{5/2} - \frac{1}{48}\log z + \sum_{g\geq2} a_g z^{-5(g-1)/2}, \quad a_g = \frac{4^{1-g}}{(3g-3)!}\ab<\tau_2^{3g-3}>_g, \]
where
\[ \ab<\tau_2^{3g-3}>_g \coloneqq \int_{\Mbar_{g,3g-3}} \psi_1^2 \cdots \psi_{3g-3}^2, \]
so we see that~\Cref{conj:cgmps} is equivalent to the statement that \(b_g = a_g\) for all \(g \geq 2\).

\subsection{The main results}%
\label{sub:The main results (Guo local P1)}

The global mirror symmetry picture of~\Cref{sec:Structure of higher-genus GW invariants of the quintic (Shuai Guo)} applies here as well. Our goal is to describe the behavior of \(F_g^{X_p}\) at the three boundary points of the one-dimensional B-model moduli space. The limit \(q \to 0\) is the large complex structure limit and gives the Gromov-Witten expansion in the A-model variable \(Q\). At \(q \to 1\), we will use the framed \(R\)-matrix to prove that \(F_g\) remains regular. The interesting point is \(q \to \frac{1}{p^2}\), which is a higher-order critical point. Unlike an ordinary conifold singularity, whose genus-\(g\) leading pole has order \(2g-2\), this degeneration has pole order at most \(5g-5\). Its leading coefficient is governed by the \((3,2)\) minimal model.

\begin{thm}%
\label{thm:local P1 mirror symmetry}
    For any integer \(p \geq 2\) and any \(g \geq 2\), the Gromov-Witten potential of \(X_p\) is computed by the Eynard-Orantin topological recursion of a spectral curve \(\mc{C}\), namely
    \[ F_g^{X_p}(Q) = (-1)^{g-1}\mc{F}_g^{X_p}(q). \]
\end{thm}

In particular, this gives a non-semiprojective toric Calabi-Yau example of the remodeling conjecture.

\begin{thm}%
\label{thm:local P1 rationality}
    For \(g \geq 2\), the potential \(F_g^{X_p}\) is the expansion of a rational function
    \[ F_g^{X_p} = (-1)^{g-1}\frac{\mc{P}_g(\Delta,p^2)}{\Delta^{5g-5}}, \qquad \Delta \coloneqq \frac{1-p^2q}{1-q} = \frac{p^2(w-w_c)}{w}, \]
    where \(\mc{P}_g \in \frac{1}{(p^2-1)^{2g-2}}\Q[p^2][\Delta]_{\leq 5g-5}\).
\end{thm}

In particular, this proves the structural ansatz~\Cref{eqn:cgmps ansatz}.

To identify the double-scaling limit, it remains to determine the constant term of \(\mc{P}_g\) at \(\Delta = 0\).

\begin{thm}%
\label{thm:local P1 constant term}
    The constant term of \(\mc{P}_g\) is
    \[ \mc{P}_g(0,p^2) = \frac{p^{6g-6}}{(p^2-1)^{2g-2}}\frac{1}{(3g-3)!}\ab<\tau_2^{3g-3}>_g. \]
\end{thm}

Because \(\wt{\mc{P}}_g(w_c,p) = \mc{P}_g(0,p^2)\ab(\frac{1}{p^2}-\frac{1}{p^4})^{5g-5}\), \Cref{thm:local P1 constant term} gives
\[ b_g = \frac{4^{1-g}}{(3g-3)!}\ab<\tau_2^{3g-3}>_g = a_g, \]
which proves~\Cref{conj:cgmps}.

\subsection{Gromov-Witten theory of \(X_p\)}%
\label{sub:Gromov-Witten theory of Xp}

For the equivariant theory of local curves and its localization formulas, see~\cite{BryanPandharipande2008}.

Because \(X_p\) is noncompact, the moduli space \(\Mbar_{g,0}(X_p,d)\) is not proper and the degree of the virtual class is not defined a priori. To define the invariants, we will introduce a torus action and pass to a twisted theory on \(\P^1\), where we can apply virtual localization.

We choose a framing \(f \in \Z\) and let \(\T_f \coloneqq \ker(\lambda_2 - f\lambda_1)\). The two localized weight rows, each of which includes one zero weight, are
\[ \begin{bmatrix} 0 & f & 1 & -f-1 \end{bmatrix}, \qquad \begin{bmatrix} -f & 0 & 1-fp+f & -1+fp \end{bmatrix}. \]
If \(f \neq 0\), then rescaling the character gives an equivalent symmetric weight vector
\[ \begin{bmatrix} 1 & -1 & -\mf{f} & \mf{f} \end{bmatrix}, \qquad \mf{f} = 1-p+\frac{2}{f}, \]
and we denote the adjusted torus by \(\T_{\mf{f}}\). If \(H^*_{\T_{\mf{f}}}(\mr{pt}) = \C[\lambda]\), this normalization is
\[ \lambda_1 = -\frac{2}{\mf{f}}\lambda, \qquad \lambda_2 = -2\lambda. \]

Let \(\pi \colon \Mbar_{g,1}(\P^1,d) \to \Mbar_{g,0}(\P^1,d)\) be the forgetful map and let \(\ev\) be the evaluation at the marking. Then we define
\[ N^{X_p}_{g,d} \coloneqq \int_{[\Mbar_{g,0}(\P^1,d)]^{\vir}_{\T_{\mf{f}}}} e_{\T_{\mf{f}}}^{-1}\ab(R^{\bullet}\pi_* \ev^*(\mc{O}(p-1)\oplus\mc{O}(-p-1))), \]
where \(e_{\T_{\mf{f}}}\) is the equivariant Euler class. For \(p=0\) this specializes to the usual resolved conifold formula.

We obtain a CohFT \(\Omega^{\lambda,\mf{f}}\) with state space the equivariant cohomology ring
\[ H^*_{\T_{\mf{f}}}(\P^1) = \frac{\C[H,\lambda]}{(H+\lambda)(H-\lambda)} \]
and pairing
\[ (a,b) \coloneqq \int_{\P^1} a \cup b \cup e^{-1}_{\T_{\mf{f}}}(\mc{O}(p-1)\oplus\mc{O}(-p-1)). \]
In the flat basis \(\ab\{\1,H\}\) the Gram matrix is given by
\[ \frac{1}{((p-\mf{f})^2-1)((p+\mf{f})^2-1)}\begin{pmatrix} 2\mf{f}p\lambda^{-3} & (1-\mf{f}^2-p^2)\lambda^{-2} \\ (1-\mf{f}^2-p^2)\lambda^{-2} & 2\mf{f}p\lambda^{-1} \end{pmatrix}. \]
We define the correlators by
\[ \Omega^{\lambda,\mf{f}}_{g,n}(\alpha_1,\ldots,\alpha_n) \coloneqq \sum_{d\geq0}\st_*\ab(\prod_{i=1}^n \ev_i^*\alpha_i \cap [\Mbar_{g,n}(\P^1,d)]^{\vir,\lambda,\mf{f}})e^{-dt}, \]
where \(\st\) is the forgetful map to \(\Mbar_{g,n}\). The twisted virtual class is given by
\[ [\Mbar_{g,n}(\P^1,d)]^{\vir,\lambda,\mf{f}} \coloneqq e^{-1}_{\T_{\mf{f}}}\ab(R^{\bullet}\pi_*\ev^*(\mc{O}(p-1)\oplus\mc{O}(-p-1))) \cap [\Mbar_{g,n}(\P^1,d)]^{\vir}_{\T_{\mf{f}}}. \]
Set \(Q = (-1)^{p+1}e^{-t}\). Then the genus-\(g\) potential is given by
\[ F^{X_p}_g(t) = \int_{\Mbar_g} \Omega^{\lambda,\mf{f}}_{g,0}. \]

\begin{rmk}
    The degree-positive invariants \(N^{X_p}_{g,d\geq1}\) are independent of \(\mf{f}\) by the topological vertex~\cite{LLLZ2009}, and their \(\lambda\)-independence also follows from~\Cref{thm:local P1 reconstruction}. We will use the \(\mf{f}\)-independence in the proof of~\Cref{thm:local P1 rationality}.
\end{rmk}

In genus zero, we have
\[ N^{X_p}_{0,d} = \frac{(-1)^{(p+1)d-1}}{d^3}\binom{p^2d-1}{d-1}, \qquad d \geq 1, \]
which allows us to compute the Yukawa coupling exactly:
\[ Y = -\partial_t^3 F_0^{X_p}(t) = \frac{1-\mf{f}^2-p^2}{((p-\mf{f})^2-1)((p+\mf{f})^2-1)} + \frac{q}{p^2q-1}. \]
Here, \(q\) and \(Q\) are related by the mirror map \(Q = q(1-q)^{p^2-1}\). Evaluating the quantum product using \((v_1 \star v_2, v_3) = \Omega^{\lambda,\mf{f}}_{0,3}(v_1,v_2,v_3)\), we obtain
\[ H \star H = -\frac{(\mf{f}^2-1)q+1}{p^2q-1}\lambda^2 \1 - \frac{2\mf{f}pq}{p^2q-1}\lambda H, \]
or in matrix form
\[ H \star = \begin{pmatrix} 0 & -\frac{(\mf{f}^2-1)q+1}{p^2q-1}\lambda^2 \\ 1 & -\frac{2\mf{f}pq}{p^2q-1}\lambda \end{pmatrix}. \]
The eigenvalues \(L_{\alpha}\) of this matrix are generically distinct. Then \(\Omega^{\lambda,\mf{f}}\) is semisimple. We write \(\ab\{e_{\alpha}\}\) for the idempotent basis, \(\Delta_{\alpha}^{-1} = (e_{\alpha},e_{\alpha})\) for the lengths, and \(\bar{e}_{\alpha} \coloneqq \Delta_{\alpha}^{1/2}e_{\alpha}\) for the normalized canonical basis.

Because the theory is semisimple, \Cref{thm:teleman} gives
\[ \Omega^{\lambda,\mf{f}} = R.T_R.\ab(\bigoplus_{\beta=0}^1 \Omega^{\mr{KW}}_{\beta}), \]
where \(\Omega^{\mr{KW}}_{\beta}\) is a one-dimensional trivial CohFT and \(T_R(z) = z(\wt{\1}-R^{-1}(z)\1)\) with \(\wt{\1} \coloneqq \sum_{\beta}\bar{e}_{\beta}\). We determine the \(R\)-matrix from the quantum differential equation
\[ z \d{S_{\tau}(z)} = \d{\tau} \star_{\tau} S_{\tau}(z). \]
Equivariant localization and quantum Riemann-Roch relate the fundamental solution to the \(R\)-matrix by
\[ (S^*_{\tau}(z)\1, \bar{e}_{\alpha})\big|_{q=0}\, C_{\alpha}(z) = (\1, R(z)\bar{e}_{\alpha})e^{u^{\alpha}(q)/z}, \]
where the initial condition \(C_{\alpha}(z)\) is an exponential involving Bernoulli numbers. Writing \(R_{i\bar{\alpha}}(z) \coloneqq (R^*(z)H^i, \bar{e}_{\alpha})\), we obtain from the quantum differential equation the second-order equation
\[ \ab( (zD+L_{\alpha})^2 + \frac{2\mf{f}pq}{p^2q-1}\lambda(zD+L_{\alpha}) + \frac{(\mf{f}^2-1)q+1}{p^2q-1}\lambda^2 )R_{0\bar{\alpha}} = 0, \]
where \(D = -\partial_t\). We can then compute the \(R\)-matrix order by order in \(z\).

\subsection{Topological recursion and the remodeling conjecture}%
\label{sub:Topological recursion and the remodeling conjecture}

We now pass to the B-model side. Our goal is to identify the reconstruction supplied by topological recursion with the Givental-Teleman reconstruction above. The Eynard--Orantin topological recursion~\cite{EO2007} takes as input a \textit{spectral curve} \((\Sigma,x,y,B)\), where \(\Sigma\) is a Riemann surface; \(x\) is locally defined with meromorphic differential \(\d{x}\) having finitely many simple zeros; \(y\) is holomorphic near those zeros, with \(\d{y} \neq 0\) there; and \(B\) is a fundamental bidifferential on \(\Sigma \times \Sigma\). Near a zero \(z^{\beta}\) of \(\d{x}\), let \(\bar{z} \neq z\) be the local conjugate point with \(x(\bar{z}) = x(z)\) and define the recursion kernel
\[ K_{\beta}(z_0,z) \coloneqq \frac{\int_{\bar{z}}^z B(z_0,z')}{2(y(z)-y(\bar{z}))\d{x}(z)}. \]
If we set \(\omega_{0,0} = \omega_{1,0} = 0\), \(\omega_{0,1}(z) = y(z)\d{x}(z)\), and \(\omega_{0,2}(z_1,z_2) = B(z_1,z_2)\), then the multi-differentials are defined recursively for \(2g-2+n+1 > 0\) by
\begin{align*}
    \omega_{g,n+1}(z_0,z_{[n]}) ={}& \sum_{\beta}\on{Res}_{z = z^{\beta}} K_{\beta}(z_0,z) \\
    &\cdot \Bigg( \omega_{g-1,n+2}(z,\bar{z},z_{[n]}) \\
    &\qquad + \sideset{}{'}\sum_{\substack{g_1+g_2 = g \\ I \sqcup J = [n]}}\omega_{g_1,\abs{I}+1}(z,z_I)\omega_{g_2,\abs{J}+1}(\bar{z},z_J)\Bigg),
\end{align*}
where the primed sum excludes \(\omega_{0,1}\). For \(g \geq 2\) the free energy is given by
\[ \mc{F}_g = \omega_{g,0} \coloneqq \frac{1}{2-2g}\sum_{\beta}\on{Res}_{z=z^{\beta}}\omega_{g,1}(z)\int_{z'=z^{\beta}}^z \omega_{0,1}(z'). \]

The remodeling conjecture~\cite{BKMP2009} of Bouchard, Klemm, Mari\~{n}o, and Pasquetti asserts that the all-genus open and closed Gromov-Witten invariants of a toric Calabi-Yau threefold \(X\) coincide with the invariants \(\omega_{g,n}\) produced by topological recursion on its mirror curve. It was proved for \(\C^3\), the resolved conifold, and local curves by Eynard-Orantin, Zhou, and Chen-Li-Liu. Then Fang--Liu--Zong~\cite{FLZ2020} proved it for all semiprojective toric Calabi-Yau threefolds. Unfortunately, \(X_p\) is not semiprojective for \(p \geq 2\), so it falls outside the scope of that proof.

To compare the two sides we use the reconstruction of Dunin-Barkowski, Orantin, Shadrin, and Spitz~\cite{DOSS2014}, which identifies the stable topological recursion correlators with the Givental reconstruction of a semisimple CohFT. Around each critical point \(z^{\beta}\), define the local Airy coordinate \(\eta^{\beta}(z)\) by
\[ x(z) = x(z^{\beta}) + \frac{1}{2}(\eta^{\beta}(z))^2 \]
and the globally defined meromorphic reference form
\[ \d{\zeta^{\bar{\beta}}}(z) \coloneqq -\on{Res}_{z'=z^{\beta}}\frac{B(z',z)}{\eta^{\beta}(z')}. \]
The state space \(V\) of the associated CohFT is spanned by the basis \(\ab\{\bar{e}_{\beta}\}\) corresponding to the critical points, and the pairing is diagonal with \(\eta(\bar{e}_{\beta},\bar{e}_{\gamma}) = \delta_{\beta,\gamma}\). The curve \(R\)-matrix \(\check{R}\) is defined by the asymptotic Laplace transform
\[ \frac{u}{\sqrt{2\pi u}}\int_{\mc{L}_{\gamma}} e^{-x(z)/u}\d{\zeta^{\bar{\beta}}}(z) \asymp e^{-x^{\gamma}/u}\eta\ab(\bar{e}_{\gamma}, \check{R}^*(-u)\bar{e}_{\beta}), \]
and the analogous transform of \(\d{y}\) defines \(T(u)\), and if the resulting CohFT has a flat unit, then \(T(u) = u(\wt{\1}-\check{R}^{-1}(u)\1)\). Then \(\check{R}\) is symplectic, the curve CohFT is \(\check{\Omega} = \check{R}.T.\ab(\bigoplus_{\beta}\Omega^{\mr{KW}}_{\beta})\), and writing \(\d{\zeta^{\bar{\beta}}_k} \coloneqq \ab(-\d{} \circ \frac{1}{\d{x}})^k \d{\zeta^{\bar{\beta}}}\), the correspondence reads
\begin{align*}
    \omega_{g,n}(z_1,\ldots,z_n) ={}& \sum_{\substack{k_1,\ldots,k_n \geq 0 \\ \beta_1,\ldots,\beta_n \in [N]}}\ab(\int_{\Mbar_{g,n}} \check{\Omega}_{g,n}(\bar{e}_{\beta_1}\psi_1^{k_1},\ldots,\bar{e}_{\beta_n}\psi_n^{k_n})) \\
    &\cdot \prod_{i=1}^n \d{\zeta^{\beta_i}_{k_i}}(z_i).
\end{align*}

For \(X_p\), we will consider the spectral curve
\[ \mc{C}_{\mf{f}} = \ab(\P^1, x_{\mf{f}}(z), y(z), B(z_1,z_2) = \frac{\d{z_1}\d{z_2}}{(z_1-z_2)^2}), \qquad x_{\mf{f}}(z) = x(z) + 2\mf{f}\lambda^2 y(z), \]
where
\begin{align*}
    x(z) &= t^0 + \lambda \log q + \lambda(p-1)\log z + \lambda(-p-1)\log\frac{q-z}{1-z}, \\
    y(z) &= \frac{p+1}{2\lambda}\log(1-q) + \frac{1}{2\lambda}\log z - \frac{1}{2\lambda}\log(1-z) - \frac{1}{2\lambda}\log(q-z).
\end{align*}
Solving \(\d{x_{\mf{f}}} = 0\), we obtain the critical points
\[ z^{\alpha} = \frac{1-\ms{L}_{\alpha}}{1+\mf{f}+p\ms{L}_{\alpha}}, \qquad \ms{L}_{\alpha} = \lambda^{-1}L_{\alpha}, \qquad q = \frac{\ms{L}_{\alpha}^2-1}{(p\ms{L}_{\alpha}+1)^2-1}, \]
and the inner products are given by \(\eta(e_{\alpha},e_{\beta}) = \delta_{\alpha,\beta}\frac{y'(z^{\alpha})^2}{x''_{\mf{f}}(z^{\alpha})}\). Then the resulting lengths satisfy \(\check{\Delta}_{\alpha} = \Delta_{\alpha}\). Therefore, the Frobenius structure of the spectral curve is the quantum product of \(X_p\).

We have now matched the Frobenius structures, so it remains to match the two \(R\)-matrices. Integration by parts on the Lefschetz thimbles gives
\begin{align*}
    u \partial_{t^0}\ab(e^{-x^{\gamma}/u}\check{R}^*(-u)^{\bar{\gamma}}_{\1}) &= -e^{-x^{\gamma}/u}\check{R}^*(-u)^{\bar{\gamma}}_H, \\
    u\partial_{t^0}\ab(e^{-x^{\gamma}/u}\check{R}^*(-u)^{\bar{\gamma}}_H) &= -\lambda b(q)e^{-x^{\gamma}/u}\check{R}^*(-u)^{\bar{\gamma}}_H - \lambda^2 c(q) e^{-x^{\gamma}/u}\check{R}^*(-u)^{\bar{\gamma}}_{\1},
\end{align*}
where
\[ c(q) = -\frac{(\mf{f}^2-1)q+1}{p^2q-1}, \qquad b(q) = -\frac{2\mf{f}pq}{p^2q-1} \]
are exactly the entries of the quantum product computed above, so these agree with the quantum differential equation for the Gromov-Witten \(R\)-matrix. To match the initial conditions, we take the limit \(q \to 0\) of the Laplace transforms along the thimbles, which produces Beta functions and therefore products of Gamma functions. Then Stirling's expansion matches the Bernoulli-number initial condition \(C_{\alpha}(z)\).

\begin{prop}%
\label{prop:R matrices coincide}
    The two \(R\)-matrices coincide, namely \(R(z) = \check{R}(u)|_{u=-z}\).
\end{prop}

We have now shown that the canonical lengths and the \(R\)-matrices agree. After the substitution \(u = -z\), the genus-\(g\) stable graph sums differ by the overall parity \((-1)^{3g-3} = (-1)^{g-1}\), which proves~\Cref{thm:local P1 mirror symmetry}.

\subsection{Rationality}%
\label{sub:Rationality}

Because the higher-genus potential is independent of \(\mf{f}\), we may, for simplicity, set \(\mf{f} = 0\). This simplifies both the canonical basis and the quantum differential equation, and the latter becomes
\[ \ab((zD+L_{\alpha})^2 - \frac{1-q}{1-p^2q}\lambda^2)R_{0\bar{\alpha}} = 0, \qquad L_{\alpha} = (-1)^{\alpha}\lambda L, \quad L = \ab(\frac{1-q}{1-p^2q})^{1/2}, \]
and we denote the entries of the \(R\)-matrix by \((R_k)_i^{\alpha}\). Set
\[ \Delta \coloneqq \frac{1}{L^2} = \frac{1-p^2q}{1-q}, \qquad \mc{R} \coloneqq \Q[p^2,\Delta], \qquad \mc{R}_{\leq d} \coloneqq \ab\{P \in \mc{R} \mid \deg_{\Delta} P \leq d\}. \]
We now use induction on the coefficients of the quantum differential equation to obtain a polynomial structure.

\begin{prop}%
\label{prop:R matrix polynomiality}
    For \(k \in \Z_{\geq0}\) and \(i = 0,1\), we have
    \[ (R_k)_i^{\alpha} \in \frac{1}{(p^2-1)^k \lambda^{k-i}\Delta^{(3k+i)/2}}\mc{R}_{\leq 2k}. \]
\end{prop}

Analyzing the degree contributions of the stable graph sum in the Givental-Teleman reconstruction, and tracking the powers of \(\lambda\), \(\Delta\), and \(p^2-1\), we obtain the following.

\begin{thm}%
\label{thm:local P1 reconstruction}
    Writing \([\,\cdot\,]_d\) for the codimension-\(d\) component, we have
    \begin{align*}
        \ab[\Omega^{\lambda}_{g,n}(H^{a_1},\ldots,H^{a_n})]_d \in{}& \frac{\lambda^{3g-3-d+\sum_i a_i}}{(p^2-1)^{d-(g-1)}\Delta^{\frac{1}{2}(g-1+3d+\sum_i a_i)}}\mc{R}_{\leq 2d} \\
        &\otimes H^{2d}(\Mbar_{g,n},\Q).
    \end{align*}
\end{thm}

Taking \(n=0\) and \(d = 3g-3 = \dim \Mbar_g\) in~\Cref{thm:local P1 reconstruction} gives
\[ F_g^{X_p} \in \frac{1}{(p^2-1)^{2g-2}\Delta^{5g-5}}\mc{R}_{\leq 6g-6}, \]
so~\Cref{thm:local P1 rationality} will follow if the degree of the numerator in \(\Delta\) is bounded by \(5g-5\). We can see this bound at \(q=1\). Because \(\Delta\) has a simple pole there, regularity of \(F_g^{X_p}\) forces the numerator to cancel the highest powers of \(\Delta\). We have therefore reduced rationality to regularity at \(q=1\). This regularity follows from the \(\mf{f}\)-independence of \(F_g^{X_p}\) together with the existence of the limit \(q \to 1\) of the \(R\)-matrix for \(\mf{f} \neq 0\).

\subsection{The limit \(\Delta = 0\)}%
\label{sub:The limit Delta = 0}

To determine the remaining constant term, we now study the limit \(\Delta = 0\), which is precisely the double-scaling limit. We rescale the canonical lengths by
\[ \delta_{\alpha} \coloneqq \Delta^{1/2}\Delta_{\alpha} = (-1)^{\alpha+1}2(p^2-1)\lambda^3 \]
and define the modified \(R\)-matrix entries by isolating the leading terms at \(\Delta = 0\) in~\Cref{prop:R matrix polynomiality}, namely
\[ (\bar{R}_k)_i^{\alpha} \coloneqq \ab[(-1)^k (R_k)_i^{\alpha}\Delta^{(3k+i)/2}]\Big|_{\Delta=0}, \qquad i = 0,1. \]
Then the graph sum for \(F_g^{X_p}\) becomes a formula for \(\mc{P}_g(0,p^2)\) after replacing \(\Delta_{\alpha}\) by \(\delta_{\alpha}\) and \((R_k)_i^{\alpha}\) by \((\bar{R}_k)_i^{\alpha}\).

\begin{lem}%
\label{lem:R matrix at Delta zero}
    Solving the leading-order quantum differential equation at \(\Delta = 0\) gives the closed forms
    \begin{align*}
        (\bar{R}_k)_0^{\alpha} &= \frac{p^{2k}}{(p^2-1)^k\lambda^k}(-1)^{k\alpha}\prod_{n=1}^k \frac{(6n-7)(6n+1)}{48n}, \\
        (\bar{R}_k)_1^{\alpha} &= \frac{p^{2k}}{(p^2-1)^k \lambda^{k-1}}(-1)^{(k-1)\alpha}\frac{1-6k}{1+6k}\prod_{n=1}^k \frac{(6n-7)(6n+1)}{48n}.
    \end{align*}
\end{lem}

After factoring out the \(p\)-dependent terms in~\Cref{lem:R matrix at Delta zero}, we obtain purely rational coefficients. We now construct the topological field theory \(\Omega^{\mr{top}}\) on \(V = \on{span}\ab\{\phi_0,\phi_1\}\) with unit \(\phi_0\), where
\[ (\phi_0,\phi_0) = (\phi_1,\phi_1) = 0, \qquad (\phi_0,\phi_1) = -1, \qquad \phi_1 \star \phi_1 = \phi_0. \]
Applying reconstruction by \(R(z)\) together with its flat-unit translation, we obtain a constant \(C_g\) with
\[ \mc{P}_g(0,p^2) = \frac{p^{6g-6}}{(p^2-1)^{2g-2}}C_g, \]
and it remains to identify the reconstructed CohFT.

The \((2,3)\) minimal model has spectral curve
\[ x(z) = \frac{1}{3}(z^3-3\ep^2 z) + \tau^0, \qquad y(z) = \frac{1}{2}(z^2-2), \]
with flat coordinate \(\tau^1 = \ep^4/2\). If we set \(\tau^0 = 0\) and \(\ep = 1\) and identify \(\phi_0 = \partial_{\tau^0}\) and \(\phi_1 = \partial_{\tau^1}\), then the \(R\)-matrix obtained from \(X_p\) at \(\Delta = 0\) coincides with the \(R\)-matrix of this curve. Therefore, \(C_g\) is exactly the free energy of the \((2,3)\) minimal model.

Finally, we apply the symplectic swap \((x,y) \mapsto (y,-x)\), under which the curve becomes
\[ x(z) = \frac{1}{2}(z^2-2), \qquad y(z) = -\frac{1}{3}(z^3-3z). \]
The swapped curve has one-dimensional state space, trivial \(R\)-matrix, and translation \(T(u) = u^2 \bar{e}\), so that
\[ C_g = \frac{1}{(3g-3)!}\ab<\tau_2^{3g-3}>_g, \]
which is~\Cref{thm:local P1 constant term}. Note that this swap is the correspondence between the \((2,3)\) and the \((3,2)\) minimal models, which is why the intersection numbers appearing here are those of~\Cref{conj:cgmps}.

